\centerline{\bf \Huge Основная теорема арифметики}
\title{Основная теорема арифметики}

\begin{flushright}
    \scriptsize Настоящая сила человека в способности измениться по своей воле.\\ Сайтама.
\end{flushright}

\textbf{\Large Основная теорема арифметики.} Любое натуральное число (кроме единицы) можно представить в виде произведения простых множителей, и притом единственным образом (с точностью до порядка сомножителей).

Такое произведение называется разложением на простые множители или каноническим разложением.

\problem{1}Найдите каноническое разложение числа 2550. Покажите, что оно делится на 6, 30, 51, 75, 85. Делится ли оно на 12, 22, 26, 27?

\problem{2}Не вычисляя произведения $2013 \cdot 15 \cdot 77$, выясните, делится ли оно на 2, 3, 9, 35, 55, 80, 6039.

\problem{3}а) Число A делится на 3 и 4. Следует ли отсюда, что A делится на 3 · 4 = 12?\\
б)Число A делится на 4 и 6. Следует ли отсюда, что A делится на 4 · 6 = 24?

\problem{4}Докажите, что произведение пяти последовательных натуральных чисел делится на 120.

\problem{5}Доказать, что из любых 27 различных натуральных чисел, меньших 100, можно выбрать два числа, не являющихся взаимно простыми.

\problem{6} На сколько нулей оканчивается число 100!?

\problem{7}а)Натуральное число умножили на каждую из его цифр. Получилось 2814. Найдите исходное число.\\
б) Натуральное число умножили на сумму его цифр и получили 1944. Найдите исходное число.

\problem{8}Расул Владимирович написал на доске $1\cdot2\cdot...\cdot100$, и сказал кто сможет это вычислить получит конфеты. Герман очень хотел конфеты и долго пытался вычислить это произведение, но у него никак это не получалось. И от отчаяния он накричал на эти числа. Но он забыл о своём стенде и числа на доске превратили в факториалы. Пытаясь вернуть всё как было, Герман решил стереть восклицательные знаки, но вместо этого стёр одно число вместе с его восклицательным знаком. Тут Баир вернулся из прошлого и глянув на доску сказал:" Да это же целый квадрат". Какое число стёр Герман? Или Баиру нужно отдохнуть? 